\subsection*{Appendix A.6: Relationship between FP4 Code Occupancy Optimization and Direct MSE Optimization}
\label{app:loss_comparison}

Compared to the approach of directly minimizing the MXFP4 quantization MSE:
\begin{equation}
\mathcal{L}_{\text{MSE}}(\mathbf{S},\mathbf{R}) = \mathbb{E}\left[ \|\mathbf{Z} - Q_{\text{mxFP4}}(\mathbf{Z})\|_F^2 \right]
\end{equation}
The proposed Codebook Occupancy Equilibrium Loss $\mathcal{L}_{\text{code}}$ offers three distinct advantages:

\begin{itemize}
    \item \textbf{Geometric Alignment:} $\mathcal{L}_{\text{code}}$ is defined entirely within the MXFP4 codeword space. It directly constrains the "usage frequency of each 4-bit codeword," aligning precisely with the "shared exponent + logarithmic quantization" structure unique to block floating-point formats\cite{mxfp}.
    \item \textbf{Optimization Smoothness:} The histogram-level loss is robust to individual outliers and is significantly smoother than point-wise MSE. In the context of coarse 4-bit quantization, $\mathcal{L}_{\text{MSE}}$ is heavily affected by the discrete quantization steps ("staircase" function), making optimization challenging due to local minima. $\mathcal{L}_{\text{code}}$ provides a smoother landscape for finding optimal rotation angles \cite{qsurvey}.
    \item \textbf{Implementation Efficiency:} The proposed method (Givens rotations + 1D angle search) requires only forward passes and counting operations. It eliminates the need for differentiable approximations, gradients, or backpropagation, making it a quintessential Post-Training Quantization (PTQ) calibration paradigm that is both fast and memory-efficient.
\end{itemize}

Empirical results demonstrate that intra-block rotation based on codeword occupancy equilibrium significantly improves MXFP4 activation quantization accuracy across various LLMs (LLaMA, Qwen). It consistently outperforms simple rotation strategies based directly on MSE in metrics such as WikiText perplexity and MMLU accuracy.